\section{Data, Taxonomy, and Measurement}
\label{sec:methods}

The article's analysis proceeds in three stages: constructing a topic taxonomy of relationship problems, generating LLM advice responses, and comparing these responses to community-endorsed advice using linguistic metrics. The taxonomy construction stage warrants methodological attention in its own right. Following Bowker and Star's framework for classification as infrastructure \citep{bowker1999sorting}, the categories developed here are treated not as neutral containers but as analytic choices that encode assumptions about where boundaries fall and what differences matter. How posts are sorted shapes what divergence becomes visible---a point that bears on the interpretation of results in Section~\ref{sec:results}.

\subsection{Corpus, Baseline, and Ethics}

Posts and comments were collected from r/relationship\_advice spanning October 2025 through January 2026, yielding 32,630 posts and 546,521 comments after quality filtering. For the primary comparison, the analysis is restricted to posts with at least one substantive, community-endorsed response: a top-level comment from someone other than the original poster with score $\geq 5$ and length $\geq 200$ characters ($n = 11{,}565$). For each eligible post, LLM-generated advice is compared against the highest-voted qualifying human comment.

R/relationship\_advice is selected as the community baseline precisely because of its structural properties, not despite them. The community satisfies the methodological conditions for defamiliarization outlined in Section~\ref{sec:bg}: concentrated, measurable norms; a distinctive vocabulary; and a prescriptive force that varies systematically across problem types. The upvote economy of Reddit selects for directiveness and moral certainty; the advice baseline constructed here likely represents a more extreme pole than typical human advice from friends, family, or strangers offline. That extremity is the analytic instrument. The goal is not to sample human advice broadly but to construct a community baseline whose norms are sufficiently concentrated and legible to make divergence from those norms specifiable. A more diffuse or ambivalent community would flatten the contrast; this one makes the comparison sharp. The community is treated as a \textit{baseline for measuring divergence}, not as a normative gold standard. Its own assumptions about exit, autonomy, and the legibility of harm are themselves subjected to critical scrutiny in Section~\ref{sec:discussion}. The divergence documented in Section 4 should therefore be read as the distance between aligned LLM defaults and one particularly concentrated normative formation, not as a deficit relative to human advice more broadly.

The highest-voted comment is not the most typical but the most rewarded under the ranking regime---a composite of content, timing, and positional effects treated here as a feature rather than a confound: the object of analysis is the normative stance that the platform's affordances select for and amplify, not the distribution of views prior to aggregation.

No attempt is made to identify individual posters; usernames are not retained and play no role in the analysis, which operates primarily at the aggregate level of linguistic metrics. Illustrative examples are drawn only from high-visibility posts (score > 500), where the assumption of public address is strongest and the advice has already entered wide community circulation.

\subsection{Constructing the Problem Taxonomy}

Standard topic models treat topics as latent distributions over vocabulary, finding lexical co-occurrence patterns rather than semantic categories \citep{Grimmer_Stewart_2013}. In this corpus, the limitation is acute: the vocabulary of relationship distress is relatively uniform across problem types. The distinction between infidelity and controlling behavior lives at the level of narrative framing, not word choice. Both LDA and BERTopic, despite parameter variation, produced only two coarse clusters, because the relevant semantic distinctions occur at the narrative level rather than the word level.

Rather than consolidating toward consensus, six models from diverse providers independently categorize posts, and their disagreement is treated as data. When six models classify the same post, convergence indicates stable conceptual boundaries; divergence indicates contested terrain where categories are not given by the facts but achieved through particular interpretive frameworks. Each model processed posts sequentially, either assigning to an existing category or proposing a new one. The 930 raw labels were consolidated into a 70-category taxonomy using Claude Sonnet 4 to merge semantically similar labels; cross-model agreement on primary topic was 17.8\%, with systematic disagreements concentrated on harm-versus-neutral framings. This contested terrain is analyzed in Section~\ref{sec:results}. Procedural details appear in Appendix~\ref{sec:supplementary}.

\subsection{Generating Model Advice}

Advice was generated using four models (Gemini 2.5 Flash Lite, GPT-4.1-nano, DeepSeek v3.2, and Ministral 8B) across all 11,565 eligible posts. These models were selected as cost-efficient, widely deployed representatives of four distinct providers and training pipelines (Google, OpenAI, DeepSeek, and Mistral), enabling comparison across alignment regimes without the prohibitive cost of frontier models. Notably, no Anthropic model was included in advice generation. Although Claude Sonnet 4 was used in the taxonomy consolidation pipeline---a post-hoc merging task whose output does not enter the comparative analysis---including it as an advice generator would create a conflict of interest in which the same model family both structured the comparison and appeared within it. It is possible that Claude-family models would show a different pattern, and the findings cannot be assumed to generalize to systems with explicitly constitutionalist alignment approaches.

Posts are presented to models with minimal scaffolding: only the post body text is provided, omitting titles and avoiding framing such as ``respond as a Reddit commenter.'' What the model chooses to do, whether advising, validating, or questioning, is itself data about its defaults. This follows work that treats prompting as an empirical audit of generative systems' defaults and normative tendencies, using systematic prompt conditions to surface patterned behavior. Temperature was set to 0.7 for advice generation (characterizing typical deployment behavior) and 0.0 for topic classification (requiring consistency). A sensitivity check at temperature 0.2 ($n = 500$) showed minimal change, with a mean absolute difference of 0.02 across key metrics and divergence persisting at comparable effect sizes ($|d| > 0.5$). This suggests the advisory stance is deeply embedded in LLMs, rather than an artifact of stochastic sampling.

\subsection{Operationalizing Advice Style}

LLM and human advice are compared using lexicon-based metrics, avoiding LLM-as-judge evaluation to prevent circularity. These measures are designed to capture robust directional differences at scale rather than fine-grained semantic content. The analysis makes claims about distributional patterns across thousands of responses, not about the deontic structure of
individual utterances. Key constructs include: \textit{epistemic stance}, operationalized through Hyland's booster/hedge distinction and the deontic/epistemic modal contrast \citep{hyland1998hedging, palmer2001mood}; \textit{relationship orientation}, captured through leave/stay language, red-flag terminology, and therapeutic vocabulary; and \textit{permission-granting}, operationalized as a four-level spectrum following Stevanovic and Per\"{a}kyl\"{a}'s framework for deontic authority transfer \citep{stevanovic2012deontic}: explicit authorization (``you can leave''), conditional authorization (``if he doesn't change, consider leaving''), values affirmation (``you deserve better''), and empathy/validation (``that sounds really hard''). The first two levels grant or approach granting permission to act; the latter two validate feelings without authorizing action. This gradient operationalization addresses construct validity concerns about binary distinctions while preserving the core finding: human advice shows a 5:1 ratio of action-oriented to validation language, while LLM advice inverts this to approximately 1:4. The analysis also measures partner-alignment prompts (``understand their perspective'') and analysis framing (``let's break down''), which index attention direction and advisory stance. Full lexicons appear in Appendix~\ref{sec:supplementary}.

These lexicons are designed as directional proxies for large-scale comparison rather than as fine-grained semantic instruments. To assess whether they recover the broad contrasts central to the argument, two independent coders compared 50 matched human--LLM response pairs on certainty, leave-orientation, and therapeutic framing. When both coders judged that a clear difference existed, they agreed on direction in 96\% of certainty comparisons (24/25 cases), 100\% of leave-orientation comparisons (20/20 cases), and 100\% of therapeutic comparisons (44/44 cases). Disagreements were primarily about whether the difference was large enough to call---one coder marking ``equal'' or ``not applicable'' where the other saw a slight advantage---rather than about which response exhibited more of the property. The near-absence of directional disagreement (only one case where coders said opposite directions) supports the paper's claim that the lexicons recover robust distributional differences in advisory style.

\subsection{Identifying High-Consensus Cases}

To test whether LLM divergence concentrates where human consensus is strongest, the analysis identifies posts where the community has reached clear agreement. Thread-level distributions are used for \textit{selection}: a post qualifies as high-consensus when at least 70\% of its qualifying human comments exhibit leave-oriented language ($n = 153$ posts). This threshold captures clear abuse patterns, discovered infidelity, and safety threats where the community's exit-orientation is most concentrated. For \textit{comparison}, consistency is maintained with the primary analysis by measuring LLM advice against the top comment baseline. On high-consensus posts, the top comment almost invariably reflects the thread's exit orientation; by definition, the voting economy surfaces the directive response when consensus exists. This design tests whether LLM caution scales with situational clarity (as calibrated systems would) or applies uniformly regardless of how certain the community is.

Robustness checks varying the consensus threshold from 50\% to 90\% show a monotonic pattern: as threshold increases, the human-LLM gap widens (0.17 at 50\%, 0.22 at 60\%, 0.30 at 70\%, 0.34 at 80--90\%). The divergence is not an artifact of threshold selection. An alternative operationalization was also considered based on score concentration (whether the top comment dominates thread visibility) but found that concentration alone does not predict divergence. The gap emerges specifically where high concentration coincides with directional consensus: visibility amplifies a settled community position that LLMs do not echo.

This design is intended to identify patterned divergence between assistant-style LLM advice and a concentrated, vote-ratified community baseline. It does not establish that the community offers better advice, nor does it isolate any single cause of model behavior. The differences observed may reflect a combination of safety alignment, training-data averaging, model architecture, provider-specific instruction tuning, and the generalized assistant stance elicited by minimal prompting. The analytic value lies elsewhere: by comparing portable model defaults with a community whose norms are unusually concentrated and legible, it becomes possible to specify what kinds of directive judgment are systematically attenuated or withheld in model outputs---and to characterize the structure of that attenuation, whatever its causal source.
